\begin{helper}
Let \(K\) be a finite field of order \(q\), and suppose \(t\ge2\).  Use the
concrete layer choice
\[
  \ell(p,r)=\noderef{StarProductLayerChoice}(p,r)
\]
and the resulting concrete marking rule
\(\noderef{StarProductConcreteMarked}\).  If \(A\) is chosen so that
\[
  10128t\le A
\]
and \(h\ge Aq^t\), then for every \(m<k\) and every legal
forward-independent prefix \(p\), the number of legal children marked by the
concrete rule is at most \(h\).
\end{helper}
\begin{proof}
Fix \(m<k\) and a legal forward-independent prefix \(p\).  The concrete
layer choice is the definition \(\noderef{StarProductLayerChoice}\) specified
in the statement.  With this choice,
\(\noderef{StarProductPopularChildrenBound}\) says that the legal popular
children of \(p\) are at most \(128tq^t\), and
\(\noderef{StarProductPoorChildrenBound}\) says that the legal poor children
of \(p\) are at most \(10000tq^t\).

The abstract counting certificate
\(\noderef{StarProductConcreteMarkedChildrenCertificate}\) applies to this
same concrete layer choice.  Its hypotheses are exactly the two concrete
popular and poor bounds just cited, together with \(10128t\le A\) and
\(Aq^t\le h\).  Therefore the set of legal children for which
\(\noderef{StarProductConcreteMarked}\) returns true has cardinality at most
\(h\), as required.
\end{proof}
